\chapter{Virtual Machines}

\begin{aside}
This module was added in Fall 2020 (listed as ``New!'' on that term's reading page). In Winter
2018 the final slot was ``Interprocess Communication and Networking'', marked as not covered, so
networking is only mentioned briefly at the end.
\end{aside}

\section{What is a virtual machine?}

\begin{definition}
A \textbf{virtual machine} (VM) is a software abstraction of a complete computer: CPU, memory,
devices. A \textbf{guest} operating system runs inside the VM. The \textbf{hypervisor} (or
\emph{virtual machine monitor}, VMM) is the software that creates and runs VMs on the real
\textbf{host} hardware, multiplexing it among the guests.
\end{definition}

Just as a process virtualizes a CPU and memory for an application, a VM virtualizes the whole
machine for an operating system.

Uses:
\begin{itemize}
  \item \textbf{Server consolidation}: many lightly loaded servers on one physical machine;
    the basis of cloud computing (AWS EC2, etc.).
  \item \textbf{Isolation}: a compromised or crashed guest does not affect others.
  \item Running different OSes side by side (development, testing, legacy software).
  \item \textbf{Snapshots and live migration}: VM state can be saved, cloned, or moved to
    another host while running.
  \item OS development and teaching --- System/161 is itself a (simulated) machine!
\end{itemize}

\section{Types of hypervisors}

\begin{center}
\begin{tikzpicture}[b/.style={draw,minimum width=1.6cm,minimum height=0.6cm,font=\small,align=center}]
  % type 1
  \node[b,minimum width=4cm,fill=gray!20] (hw1) at (0,0) {hardware};
  \node[b,minimum width=4cm,fill=blue!10,above=0pt of hw1] (hv1) {hypervisor};
  \node[b,above=0pt of hv1.north west,anchor=south west,minimum width=2cm] (g1) {guest OS};
  \node[b,above=0pt of hv1.north east,anchor=south east,minimum width=2cm] (g2) {guest OS};
  \node[b,above=0pt of g1,minimum width=2cm] {apps};
  \node[b,above=0pt of g2,minimum width=2cm] {apps};
  \node[below=0.1cm of hw1,font=\small] {Type 1 (bare metal)};
  % type 2
  \node[b,minimum width=4cm,fill=gray!20] (hw2) at (6,0) {hardware};
  \node[b,minimum width=4cm,above=0pt of hw2] (host) {host OS};
  \node[b,above=0pt of host.north west,anchor=south west,minimum width=2cm,fill=blue!10] (hv2) {hypervisor};
  \node[b,above=0pt of host.north east,anchor=south east,minimum width=2cm] (ha) {host apps};
  \node[b,above=0pt of hv2,minimum width=2cm] (go) {guest OS};
  \node[b,above=0pt of go,minimum width=2cm] {apps};
  \node[below=0.1cm of hw2,font=\small] {Type 2 (hosted)};
\end{tikzpicture}
\end{center}

\begin{description}
  \item[Type 1 (native, bare metal)] runs directly on the hardware: Xen, VMware ESXi,
    Microsoft Hyper-V. Often a small privileged guest (Xen's \emph{dom0}) handles device drivers.
  \item[Type 2 (hosted)] runs as an application on a host OS and uses the host's drivers:
    VMware Workstation, VirtualBox, QEMU. Linux KVM blurs the line (the Linux kernel itself
    becomes the hypervisor).
\end{description}

\section{Virtualizing the CPU}

\subsection{Trap and emulate}
Run the guest OS \emph{directly} on the CPU, but in user mode (deprivileged). Ordinary
instructions run at full speed. When the guest executes a privileged instruction (e.g.\ write the
TLB, disable interrupts, return from exception), the CPU traps to the hypervisor, which
\textbf{emulates} the effect on the VM's \emph{virtual} state (virtual CPU registers, virtual
interrupt flag, virtual TLB) and returns.

System calls and exceptions from guest applications also trap into the hypervisor first; it
reflects them into the guest OS by jumping to the guest's exception handler as if the hardware
had done so (and records that the virtual CPU is now in ``virtual kernel mode'').

\begin{figure}[H]
\centering
\begin{tikzpicture}[x=0.85cm, y=1cm, font=\small, >=Stealth]
  \node[anchor=east] at (-0.1,1.6) {guest kernel};
  \node[anchor=east, align=right] at (-0.1,0.8) {(real user mode)};
  \node[anchor=east] at (-0.1,0) {hypervisor};
  \fill[blue!6] (0,-0.4) rectangle (11,0.4);
  \draw[line width=4pt, green!50!black!60] (0,1.6) -- (2.5,1.6);
  \node[font=\scriptsize] at (1.25,1.95) {ordinary instructions};
  \node[font=\scriptsize] at (1.25,1.25) {(full speed)};
  \draw[->, thick, red!70!black] (2.5,1.5) -- node[right, font=\scriptsize]{privileged instr.\ traps} (2.8,0.1);
  \draw[line width=4pt, orange!70] (2.8,0) -- (6.2,0);
  \node[font=\scriptsize, align=center] at (4.5,-0.85) {emulate on the VM's \emph{virtual} state\\(virtual TLB, interrupt flag, \ldots)};
  \draw[->, thick] (6.2,0.1) -- (6.5,1.5);
  \draw[line width=4pt, green!50!black!60] (6.5,1.6) -- (11,1.6);
  \node[font=\scriptsize] at (8.7,1.95) {continue after the instruction};
\end{tikzpicture}
\caption{Trap and emulate: the deprivileged guest kernel runs directly on
the CPU until it executes a privileged instruction, which traps to the
hypervisor.}
\end{figure}

\begin{theorem}[Popek and Goldberg, 1974, informally]
A machine can be virtualized by trap-and-emulate if every \emph{sensitive} instruction (one that
reads or changes privileged machine state or behaves differently in user vs.\ kernel mode) is
also \emph{privileged} (traps when executed in user mode).
\end{theorem}

Classic x86 violated this: e.g.\ \code{popf} in user mode silently ignores changes to the
interrupt flag instead of trapping. Solutions:
\begin{itemize}
  \item \textbf{Binary translation} (VMware, 1999): scan guest kernel code before running it and
    rewrite problematic instructions into calls to the hypervisor; cache translated blocks.
  \item \textbf{Paravirtualization} (Xen): modify the guest OS to call the hypervisor explicitly
    (\textbf{hypercalls}) instead of using sensitive instructions. Fast, but requires guest source
    changes.
  \item \textbf{Hardware-assisted virtualization} (Intel VT-x, AMD-V, 2005+): a new CPU mode
    (``guest mode'' / non-root mode) in which the guest kernel runs at its normal privilege
    level; specified events cause a \emph{VM exit} to the hypervisor, which later does a
    \emph{VM entry}. Guest state is kept in a control structure (VMCS).
\end{itemize}

\section{Virtualizing memory}

Now there are three kinds of addresses:
\[
  \text{guest virtual} \xrightarrow{\text{guest page table}} \text{guest physical}
  \xrightarrow{\text{hypervisor}} \text{host (machine) physical}.
\]
The guest OS believes it controls physical memory; the hypervisor actually allocates machine
frames to the VM.

\begin{figure}[H]
\centering
\begin{tikzpicture}[font=\small, >=Stealth, thick,
  as/.style={draw, rounded corners, minimum width=2.2cm, minimum height=0.9cm, align=center}]
  \node[as, fill=blue!10] (gv) at (0,0) {guest virtual};
  \node[as, fill=green!12] (gp) at (5,0) {guest physical};
  \node[as, fill=orange!15] (mp) at (10,0) {machine physical};
  \draw[->] (gv) -- node[above, font=\scriptsize]{guest page table} (gp);
  \draw[->] (gp) -- node[above, font=\scriptsize]{hypervisor's map} (mp);
  \draw[->, red!70!black, densely dashed] (gv.south) to[out=-25, in=-155] node[below, font=\scriptsize]{shadow page table (used by the real MMU)} (mp.south);
\end{tikzpicture}
\caption{Two levels of translation under virtualization. A shadow page
table composes them so the hardware sees a single mapping; nested page
tables (EPT/NPT) let the hardware walk both instead.}
\end{figure}

\begin{description}
  \item[Shadow page tables] The hypervisor write-protects the guest's page tables; whenever
    the guest changes them, it traps, and the hypervisor updates a \emph{shadow} page table that
    maps guest virtual directly to machine physical. The real MMU uses the shadow table. Correct
    but many expensive traps.
  \item[Software-managed TLB (MIPS)] On a TLB miss the hypervisor's handler gets control, runs the
    guest's refill handler (whose \code{tlbwr} traps), translates the guest physical frame to the
    machine frame, and installs the combined mapping.
  \item[Nested / extended page tables] (Intel EPT, AMD NPT) The hardware walks \emph{both}
    tables: guest page table (guest virtual $\to$ guest physical) and the hypervisor's table
    (guest physical $\to$ machine). No traps on guest page table updates, but a TLB miss is
    expensive: with two 4-level tables a full walk can take up to $24$ memory references
    ($(4+1)(4+1)-1$), so large pages and big TLBs matter.
\end{description}

Memory overcommitment: the hypervisor can page guest memory to disk, use a \textbf{balloon
driver} inside the guest (the balloon ``allocates'' guest pages so the guest's own replacement
policy chooses what to give up, and the hypervisor reclaims the backing frames), and share
identical pages between VMs copy-on-write (\textbf{content-based page sharing}).

\section{Virtualizing I/O}

\begin{itemize}
  \item \textbf{Emulation}: present a well-known virtual device (e.g.\ an old network card or IDE
    disk); the guest's ordinary driver accesses its registers, each access traps, and the
    hypervisor emulates the device using real host devices. Works with unmodified guests but slow.
  \item \textbf{Paravirtual devices} (virtio, Xen split drivers): the guest uses a special driver
    that communicates with the hypervisor through shared-memory ring buffers, batching requests
    and avoiding per-register traps.
  \item \textbf{Device passthrough}: give a VM direct access to a physical device (or a virtual
    function of an SR-IOV device), with an IOMMU translating and restricting the device's DMA
    addresses so it can only reach that VM's memory.
\end{itemize}

\section{Containers (OS-level virtualization)}
Containers (Docker, LXC) isolate groups of processes that \emph{share one kernel}: namespaces give
each container its own view of PIDs, file system mounts, network interfaces and users; cgroups
limit their resource usage. Much lighter than VMs (no guest kernel), but weaker isolation and the
same OS kernel for everyone.

\begin{figure}[H]
\centering
\begin{tikzpicture}[b/.style={draw, minimum height=0.6cm, font=\small, align=center}]
  \node[b, minimum width=4.4cm, fill=gray!20] (hw1) at (0,0) {hardware};
  \node[b, minimum width=4.4cm, fill=blue!10, above=0pt of hw1] (hv) {hypervisor};
  \foreach \x in {-1.47,0,1.47} {
    \node[b, minimum width=1.47cm, fill=orange!15] at (\x,1.2) {\scriptsize guest OS};
    \node[b, minimum width=1.47cm] at (\x,1.8) {\scriptsize app};
  }
  \node[font=\small] at (0,-0.65) {virtual machines};
  \begin{scope}[xshift=6.2cm]
    \node[b, minimum width=4.4cm, fill=gray!20] (hw2) at (0,0) {hardware};
    \node[b, minimum width=4.4cm, fill=orange!15, above=0pt of hw2] (k) {one shared kernel};
    \foreach \x in {-1.47,0,1.47} {
      \node[b, minimum width=1.47cm, fill=green!10] at (\x,1.2) {\scriptsize container};
      \node[b, minimum width=1.47cm] at (\x,1.8) {\scriptsize app};
    }
    \node[font=\small] at (0,-0.65) {containers};
  \end{scope}
\end{tikzpicture}
\caption{Each VM runs its own guest kernel; containers are isolated groups
of processes on a single shared kernel.}
\end{figure}

\section{Networking (brief, for context)}
Not examined in the terms covered here, but for completeness: the OS implements the network
protocol stack (link layer driver, IP routing, TCP/UDP) in the kernel and exposes it via the
\textbf{socket} API (\code{socket}, \code{bind}, \code{listen}, \code{accept}, \code{connect},
\code{send}, \code{recv}). Sockets are also the main mechanism for inter-process communication
between machines, complementing pipes and shared memory on a single machine.

\begin{lstlisting}[style=c]
int s = socket(AF_INET, SOCK_STREAM, 0);          /* TCP socket */
struct sockaddr_in addr = { .sin_family = AF_INET,
                            .sin_port = htons(8080),
                            .sin_addr.s_addr = htonl(INADDR_ANY) };
bind(s, (struct sockaddr *)&addr, sizeof addr);
listen(s, 16);
int c = accept(s, NULL, NULL);                    /* blocks for a client */
char buf[256]; ssize_t n = recv(c, buf, sizeof buf, 0);
send(c, buf, n, 0);                               /* echo it back */
close(c); close(s);
\end{lstlisting}

\section*{Summary}
\begin{itemize}
  \item A hypervisor multiplexes hardware among guest OSes; Type 1 runs on bare metal, Type 2 on
    a host OS.
  \item CPU: trap-and-emulate works when all sensitive instructions are privileged
    (Popek--Goldberg); otherwise binary translation, paravirtualization, or hardware support
    (VT-x/AMD-V).
  \item Memory: guest virtual $\to$ guest physical $\to$ machine physical, via shadow page
    tables or hardware nested page tables; ballooning and page sharing for overcommit.
  \item I/O: device emulation, paravirtual drivers, or passthrough with an IOMMU.
  \item Containers share the host kernel and are lighter but less isolated than VMs.
\end{itemize}
